%% kapa: exam=1 round=37 year=2026 period=1 form=A subject=civil_law subject_ko=민법 seq=22 number=22
%% source: Q-Net 감정평가사 기출문제 (cst00309) artlSeq=5258924 file=raw/1008/5258924/2026년 제37회 감정평가사 1차 1교시.pdf
%% conversion: kordoc 4.24.51 -> tools/gen1.py (auto); review=proofread-sonnet; flags=none
\begin{question}{22}
물권적 청구권에 관한 설명으로 옳지 않은 것은? (다툼이 있으면 판례에 따름)
\begin{choices}
  \item 소유권을 상실한 전(前)소유자는 물권적 청구권을 행사할 수 없다.
  \item 토지소유자는 그 소유 토지 위에 무단으로 건축된 건물을 임차하여 점유하고 있는 자를 상대로 건물철거를 청구할 수 있다.
  \item 지상권을 설정한 토지소유자도 불법점유자에 대하여 물권적 청구권을 행사할 수 있다.
  \item 미등기건물을 매수한 사람은 소유권이전등기를 갖출 때까지 그 건물의 불법점유자에게 직접 자신의 소유권에 기하여 인도를 청구하지 못한다.
  \item 민법은 물권적 청구권을 실현시키기 위한 비용부담의 주체에 관하여 규정하고 있지 않다.
\end{choices}
\end{question}
\answer{2}
